\section{Scope, physical coordinates, and applications}\label{b:comparison}
Theorem~\ref{b:approximation} uses a spectral bound, rather than a
small operator norm. For example,
\[
 A=\begin{pmatrix}0&2\\1/8&0\end{pmatrix}
\]
has eigenvalues $\pm1/2$ and Euclidean operator norm $2$.
Its linear field satisfies the theorem with $L=1$: one off-diagonal
coefficient is large, but their product is below one. A square-zero
matrix can have arbitrarily large norm and still satisfies the theorem
for every fixed segment length. In either case the inverse bound in
Section~\ref{b:linear-coverings} is finite for the fixed field and length;
no uniform number of cylinders is asserted as the spectral margin tends
to zero or the derivative norm grows.

Physical rescaling must be kept separate from rotating a tile.
The map $T_h=\operatorname{diag}(1,h)$ conjugates a derivative to
$T_hAT_h^{-1}$ and thus preserves its trace, determinant, eigenvalues
and square-zero property. It need not preserve the ordinary operator
norm. The weight in Proposition~\ref{b:projection} is always $hJ_h(g)$,
with $g$ expressed in the original label coordinates. This gives the
same physical cost whether the tile was selected in the original axes
or in a locally chosen orthonormal basis.

The square construction handles arbitrary compact label sets by
integration over outer neighborhoods. The product and logarithmic
constructions instead use label domains with area-zero boundary so that
almost all label area can be packed by interior cells. Their remaining
labels are covered by inexpensive open cylinders. Those cell bases are
bounded and open and need not be convex; the principal construction
produces compact parallelogram perpendicular bases. All costs count
multiplicity, including overlaps between different coarse patches or
between cell and residual cylinders.

Finally, the field, physical scale, and axial interval are fixed before
the tolerance and mesh choices. In an application with a singular
endpoint, first restrict to a compact label set contained in the field's
differentiability domain. The omitted segments, including the endpoint
itself, require a separate cover. Once the integral cost plus that extra
cost has a strict saving, the finite approximation error can be chosen
within the remaining gap. No mesh uniform over the later parameter limit
is needed.
